\documentclass[aspectratio=169]{beamer}
\usepackage[utf8]{inputenc}
\usepackage[spanish]{babel}

% Paquete oficial de plantilla Beamer de Icesi
\usepackage{custombeamer}

\begin{document}

% ==============================================================================
% PORTADA Y AGENDA
% ==============================================================================
\titleSlideSplit
  {DevOps \& CI/CD}
  {Guía Integral de Arquitectura: Scripts, Docker, GitFlows y GitHub Actions}
  {Departamento de TIC · Ingeniería de Software}
  {Universidad Icesi · 2026}

\agendaSlideCards
  {Ruta de Aprendizaje}
  {Módulo 1: Scripts de Automatización}
  {Módulo 2: Contenedores con Docker}
  {Módulo 3: Flujos de Git (GitFlow)}
  {Módulo 4: CI/CD con GitHub Actions}

% ==============================================================================
% MÓDULO 1: AUTOMATIZACIÓN CON SCRIPTS (BASH & SHELL SCRIPTING)
% ==============================================================================
\sectionSlideStripe
  {01}
  {Scripts de Automatización}
  {Principios de ejecución defensiva, validación de parámetros y robustez}

\contentSlideBulletCard
  {Shell Scripting Defensivo}
  {
    \begin{itemize}
      \item \textbf{set -e:} Detiene el script ante el primer comando con código de salida no cero.
      \item \textbf{set -u:} Lanza error si se intenta expandir una variable no inicializada.
      \item \textbf{set -o pipefail:} Propaga códigos de error en tuberías (\texttt{cmd1 | cmd2}).
      \item \textbf{trap on\_exit:} Limpieza automática de archivos temporales al terminar.
    \end{itemize}
  }
  {Principio de Robustez}
  {En servidores y pipelines automatizados, un script jamás debe continuar ejecutándose ante fallos silenciosos.}

\contentSlideTwoCols
  {Validación de Argumentos y Entorno}
  {Parámetros y Variables}
  {
    \begin{itemize}
      \item Validación de argumentos posicionales obligatorios (\texttt{\$1}, \texttt{\$2}).
      \item Asignación de valores por defecto: \texttt{\$\{APP\_ENV:-dev\}}.
      \item Carga segura de variables de entorno desde archivos \texttt{.env}.
    \end{itemize}
  }
  {Verificaciones Previas}
  {
    \begin{itemize}
      \item Verificación de herramientas instaladas con \texttt{command -v}.
      \item Comprobación de disponibilidad de puertos y bases de datos.
      \item Validación de permisos y espacio en disco disponible.
    \end{itemize}
  }

\contentSlideTwoCols
  {Trazabilidad y Manejo de Errores}
  {Códigos de Salida}
  {
    \begin{itemize}
      \item \textbf{exit 0:} Ejecución exitosa y condiciones cumplidas.
      \item \textbf{exit 1..255:} Códigos de error específicos para depuración.
      \item Uso de constantes con nombres descriptivos de errores.
    \end{itemize}
  }
  {Formato de Mensajes}
  {
    \begin{itemize}
      \item Diferenciación entre mensajes informativos y de error (\texttt{stderr}).
      \item Salidas con marcas de tiempo y niveles (\texttt{INFO}, \texttt{WARN}, \texttt{ERROR}).
      \item Formato legible y compatible con sistemas de agregación de logs.
    \end{itemize}
  }

\contentSlideQuoteStat
  {Automatización Confiable}
  {set -euo pipefail}
  {Estándar de la Industria}
  {“La automatización defensiva previene el error humano y garantiza ejecuciones determinísticas en cualquier infraestructura.”}

% ==============================================================================
% MÓDULO 2: CONTENERIZACIÓN CON DOCKER & DOCKER COMPOSE
% ==============================================================================
\sectionSlideSolid
  {02}
  {Contenerización con Docker}
  {Multi-Stage Builds para Spring Boot y React, Compose y Reverse Proxy}

\contentSlideBulletCard
  {Multi-Stage Build: Backend Spring Boot}
  {
    \begin{itemize}
      \item \textbf{Etapa 1 (Build):} Imagen JDK completa con Gradle o Maven para compilar el archivo JAR.
      \item \textbf{Etapa 2 (Runtime):} Imagen JRE mínima de ejecución (Eclipse Temurin Alpine).
      \item \textbf{Seguridad:} Creación y uso de un usuario no privilegiado (\texttt{USER appuser}).
      \item \textbf{Resultado:} Imagen final ligera (< 150MB) sin compiladores ni código fuente.
    \end{itemize}
  }
  {Aislamiento de Construcción}
  {El código fuente y las herramientas de compilación quedan aislados de la imagen final desplegada en producción.}

\contentSlideBulletCard
  {Multi-Stage Build: Frontend React SPA}
  {
    \begin{itemize}
      \item \textbf{Etapa 1 (Build):} Contenedor Node.js para instalar dependencias (\texttt{npm ci}) y empaquetar con Vite.
      \item \textbf{Etapa 2 (Runtime):} Servidor web Nginx Alpine ultraligero (< 25MB).
      \item \textbf{Configuración SPA:} Redirección de rutas HTML5 (\texttt{try\_files \$uri /index.html}).
      \item \textbf{Rendimiento:} Servido de estáticos optimizado con compresión gzip habilitada.
    \end{itemize}
  }
  {Eficiencia de Despliegue}
  {La imagen final no incluye Node.js ni la carpeta \texttt{node\_modules}, reduciendo drásticamente la superficie de ataque.}

\contentSlideTwoCols
  {Orquestación con Docker Compose}
  {Definición de la Arquitectura}
  {
    \begin{itemize}
      \item \textbf{Servicio DB:} PostgreSQL con volúmenes persistentes nombrados.
      \item \textbf{Servicio API:} Spring Boot conectado a la BD mediante variables \texttt{.env}.
      \item \textbf{Servicio Web:} Frontend React sirviendo la interfaz de usuario.
    \end{itemize}
  }
  {Aislamiento de Redes}
  {
    \begin{itemize}
      \item Creación de una red bridge interna para comunicación inter-servicios.
      \item Resolución DNS interna por nombre de servicio (e.g. \texttt{postgres-db:5432}).
      \item Mapeo exclusivo de puertos públicos necesarios hacia el exterior.
    \end{itemize}
  }

\contentSlideTwoCols
  {Enrutamiento con Nginx Reverse Proxy}
  {Punto de Entrada Unificado}
  {
    \begin{itemize}
      \item Recepción centralizada de tráfico HTTP/HTTPS en el host.
      \item Terminación SSL/TLS segura y gestión de certificados.
      \item Enrutamiento por prefijo de ruta o subdominio.
    \end{itemize}
  }
  {Mapeo de Rutas de la Aplicación}
  {
    \begin{itemize}
      \item \texttt{/}: Reenvío al contenedor del Frontend React (puerto 80).
      \item \texttt{/api/}: Reenvío al contenedor de la API Spring Boot (puerto 8080).
      \item Inyección de cabeceras \texttt{Host}, \texttt{X-Real-IP} y \texttt{X-Forwarded-For}.
    \end{itemize}
  }

\contentSlideQuoteStat
  {Arquitectura en Contenedores}
  {Microservicios}
  {Portabilidad y Aislamiento}
  {“Contenerizar cada capa de la aplicación desacopla el desarrollo de la infraestructura y simplifica la escalabilidad.”}

% ==============================================================================
% MÓDULO 3: FLUJOS DE TRABAJO EN GIT (GITFLOW) & VERSIONADO
% ==============================================================================
\sectionSlideStripe
  {03}
  {Estrategias de Git y GitFlow}
  {Topología de ramas, versionado semántico y control de calidad}

\contentSlideTwoCols
  {Estructura de Ramas en GitFlow}
  {Ramas Permanentes}
  {
    \begin{itemize}
      \item \textbf{main / master:} Representa el estado en producción. Contiene solo releases estables y etiquetados.
      \item \textbf{develop:} Rama de integración continua donde convergen las nuevas funcionalidades terminadas.
    \end{itemize}
  }
  {Ramas de Soporte Temporal}
  {
    \begin{itemize}
      \item \textbf{feature/*:} Desarrollos individuales nacidos de \texttt{develop}.
      \item \textbf{release/*:} Preparación y estabilización de versiones.
      \item \textbf{hotfix/*:} Corrección urgente de bugs críticos en producción.
    \end{itemize}
  }

\contentSlideBulletCard
  {Versionado Semántico (SemVer 2.0.0)}
  {
    \begin{itemize}
      \item \textbf{MAJOR (X.0.0):} Cambios incompatibles que rompen contratos de API o modelos de datos.
      \item \textbf{MINOR (1.X.0):} Incorporación de nuevas funcionalidades retrocompatibles.
      \item \textbf{PATCH (1.0.X):} Corrección de errores y bugs sin alterar el comportamiento general.
      \item \textbf{Sufijos:} \texttt{-SNAPSHOT} para ramas de desarrollo y \texttt{-RC1} para candidatos a release.
    \end{itemize}
  }
  {Contrato de Versiones}
  {SemVer permite a los clientes y consumidores de un API predecir el impacto de actualizar una dependencia o servicio.}

\contentSlideTwoCols
  {Commits Convencionales y Trazabilidad}
  {Tipos de Commit Estandarizados}
  {
    \begin{itemize}
      \item \textbf{feat:} Nueva funcionalidad en backend o frontend.
      \item \textbf{fix:} Corrección de un fallo o error reportado.
      \item \textbf{refactor:} Mejora estructural del código sin cambiar comportamiento.
      \item \textbf{ci / docs / test:} Tareas de soporte y pipeline.
    \end{itemize}
  }
  {Estructura del Mensaje}
  {
    \begin{itemize}
      \item Encabezado conciso en modo imperativo con alcance opcional.
      \item Cuerpo descriptivo explicando el \textit{por qué} del cambio.
      \item Referencias a tickets, issues o cambios disruptivos (\texttt{BREAKING CHANGE}).
    \end{itemize}
  }

\contentSlideBulletCard
  {Control de Calidad: Pull Requests y Code Review}
  {
    \begin{itemize}
      \item \textbf{Protección de Ramas:} Restringir el push directo a \texttt{main} y \texttt{develop}.
      \item \textbf{Pruebas Automatizadas:} Ejecución obligatoria del pipeline de CI antes de autorizar el merge.
      \item \textbf{Revisión por Pares:} Aprobación requerida de al menos un revisor del equipo de desarrollo.
      \item \textbf{Estrategia de Merge:} Uso de \textit{Squash and Merge} para mantener un historial limpio y lineal.
    \end{itemize}
  }
  {Cultura de Calidad}
  {El Pull Request es el punto focal de colaboración técnica, detección temprana de deuda técnica y transferencia de conocimiento.}

\contentSlideQuoteStat
  {Integridad del Código}
  {Branch Protection}
  {Disciplina en Repositorios}
  {“Ningún cambio ingresa a las ramas principales sin validación automatizada de pruebas y revisión de código.”}

% ==============================================================================
% MÓDULO 4: CI/CD CON GITHUB ACTIONS
% ==============================================================================
\sectionSlideDark
  {04}
  {CI/CD con GitHub Actions}
  {Pipelines declarativos, workflows condicionales y entornos de ejecución}

\contentSlideBulletCard
  {Anatomía de un Pipeline Declarativo}
  {
    \begin{itemize}
      \item \textbf{Ubicación:} Archivos YAML estructurados bajo \texttt{.github/workflows/}.
      \item \textbf{Triggers (on):} Eventos de activación por \texttt{push}, \texttt{pull\_request} o manuales (\texttt{workflow\_dispatch}).
      \item \textbf{Jobs:} Conjunto de tareas que se ejecutan en paralelo o de forma secuencial con dependencias (\texttt{needs}).
      \item \textbf{Steps:} Pasos individuales que invocan acciones de la comunidad (\texttt{uses}) o comandos de shell (\texttt{run}).
    \end{itemize}
  }
  {Pipeline as Code}
  {Toda la lógica de integración, pruebas y despliegue vive versionada en el mismo repositorio del proyecto.}

\contentSlideTwoCols
  {Pipeline de Integración Continua (CI)}
  {Fase de Backend (Spring Boot)}
  {
    \begin{itemize}
      \item Configuración del JDK mediante \texttt{actions/setup-java}.
      \item Ejecución de pruebas unitarias y de integración (\texttt{mvn test} / \texttt{gradle test}).
      \item Análisis estático de código y reporte de cobertura de pruebas.
    \end{itemize}
  }
  {Fase de Frontend (React)}
  {
    \begin{itemize}
      \item Instalación determinística con \texttt{npm ci}.
      \item Validación de sintaxis y estilos con Linters (\texttt{ESLint}).
      \item Ejecución de pruebas unitarias de componentes con \texttt{Jest} o \texttt{Vitest}.
    \end{itemize}
  }

\contentSlideTwoCols
  {Estrategia de Despliegues Condicionales}
  {Validación en Pull Requests}
  {
    \begin{itemize}
      \item Se activa ante \texttt{pull\_request} hacia \texttt{develop} o \texttt{main}.
      \item Ejecuta compilación y pruebas sin realizar despliegues en servidores.
      \item Bloquea la integración si alguna prueba unitaria falla.
    \end{itemize}
  }
  {Despliegue Continuo (CD)}
  {
    \begin{itemize}
      \item Se activa únicamente tras el merge exitoso a \texttt{main} o \texttt{release/*}.
      \item Construye las imágenes Docker optimizadas de la API y el Frontend.
      \item Despliega los nuevos contenedores y valida la salud del servicio.
    \end{itemize}
  }

\contentSlideBulletCard
  {Entornos de Ejecución: Cloud vs Self-Hosted Runners}
  {
    \begin{itemize}
      \item \textbf{GitHub-Hosted Runners:} Máquinas virtuales efímeras en la nube pública; ideales para pruebas unitarias aisladas.
      \item \textbf{Self-Hosted Runners:} Agentes ejecutándose en servidores propios o institucionales de la universidad.
      \item \textbf{Seguridad de Red:} Conexión por sondeo seguro (\textit{long-polling}) sin exponer puertos públicos de entrada.
      \item \textbf{Acceso Interno:} Despliegue directo sobre redes privadas, sockets de Docker locales y bases de datos internas.
    \end{itemize}
  }
  {Flexibilidad Híbrida}
  {Se pueden combinar jobs en la nube para análisis de código y jobs en runners locales para despliegue en servidores del laboratorio.}

\contentSlideQuoteStat
  {Entrega Continua}
  {CI / CD}
  {Ciclo de Vida Automatizado}
  {“Un pipeline robusto transforma el código fuente en un servicio en producción de forma rápida, repetible y confiable.”}

% ==============================================================================
% CIERRE Y RECAPITULACIÓN
% ==============================================================================
\agendaSlideProgress
  {Checklist DevOps para Estudiantes}
  {4}
  {1. Scripts con set -euo pipefail y validación}
  {2. Dockerfiles Multi-Stage para Spring y React}
  {3. Disciplina en Ramas de GitFlow y SemVer}
  {4. Pipelines Declarativos con GitHub Actions}

\end{document}
